\scanpage{196}
% archive sheet number: 195
% source lines 1--5
Cette \uline{filtration est fonctorielle} $\unsure{\underline{C}|\underline{\mathcal{C}}}$ loc. de longueur finie $\Leftrightarrow$ tout $X$ est le $\sup$ des $F_i(X)$. Soit en général $\Phi_0(X)=\operatorname{Sup}F_i(X)$, c'est le plus grand sous-objet de $X$ qui est ``de dim de Krull 0''.
% source insertion after line 5, continued between the lines and in the left margin
\inserted{C'est le noyau de $X\longrightarrow X_{C_0}$ (localisé par $C_0$) où $X_{C_0}$ est le \unsure{localisé} de $X$ pour la \unsure{ss-catégorie} $C_0$ de $C$ des objets de dim de Krull 0}
% slanted continuation of the insertion: de dim de Krull 0
% source lines 6--11
\unsure{On} introduit de \unsure{\deleted{même}|même} $\Phi_1(X)$ comme l'image inverse de $\Phi_0(X_{C_0}\deleted{\unsure{/\Phi_0(X)}})$\inserted{ou comme $\operatorname{Ker}(X\longrightarrow X_{C_1})$}. On en \unsure{trouve} encore une \uline{filtration} \uline{fonctorielle} $(\Phi_i(X))_i$. Dire que $\underline{C}$ est de dim de Krull $\leq n$ signifie que $\Phi_n(X)=X$ pour tout $X$.
